\section{验证器监督的回看}
\label{sec:method}
两个实现都从同一份验证器输出中学习"该回读哪一屏"。它们的差别在于历史如何抵达模型，这也意味着同样的 188 个样例对其中一个产生 2,861 条训练项、对另一个只产生 188 条。

\paragraph{标签与样本对。} 每个已验证样例 $(e,t)$ 对每个更早屏 $i<t$ 产生一条训练项。对逐屏打分器，这条训练项是一个样本对：输入包含任务描述、第 $t$ 屏的 OCR 文本、第 $i$ 屏的 OCR 文本以及距离 $t-i$；若 $i$ 在验证器给出的集合内则标签为 1，否则为 0。188 个样例共产生 2,861 个样本对，其中 463 个为正例。对全历史读取器，每个样例只产生一条训练项：输入包含任务与每个更早屏的 OCR 文本，目标是验证器集合中最晚那一屏的屏号。没有任何训练项由人标注。

\paragraph{逐屏打分器。} Qwen2.5-7B-Instruct~\citep{qwen25}，在注意力投影上加秩为 16 的 LoRA（低秩适配）适配器（4,000 万可训练参数）。prompt 以"Answer Yes or No"结尾；我们取答案位置上 \texttt{Yes} 与 \texttt{No} 的 logits，做二元 softmax，并以正例权重 3 比 1 最小化交叉熵。定位取各更早屏 $P(\text{Yes})$ 的最大者，由 $t$ 次每次 600--900 token 的前向计算得到，不做生成。

\paragraph{全历史读取器。} 同一基座模型与适配器秩，微调为从包含整段历史的 prompt 中输出屏号，损失只施加在屏号上。定位是一次最多 5 个 token 的贪心生成。两个变体都训练 2 个 epoch，学习率 $10^{-4}$，每次更新 4 条训练项，在一张 48GB GPU 上每折 13--27 分钟。

\paragraph{开销。} 在 $t=16$ 时，打分器分 $t$ 次小的前向共读入约 12K token；全历史读取器一次读入 1--5K token，prompt 全历史与之相同。3B 规模的打分器整体下降 4.2 个点，但在 $t\ge20$ 上提升 11.1 个点（\S\ref{sec:ablation}），因此打分器的开销是可调的。

\paragraph{读取阶段。} 定位之后，由一个 Qwen2.5-14B-Instruct 读取器输出要输入的文字；所有条件共用同一 prompt。\textbf{单屏读取}收到任务、当前屏以及选中屏的 OCR 文本。\textbf{带历史读取}额外收到完整历史，并把选中屏重复一遍、标注后置于当前屏之前；其开销是全历史的上下文再加一屏的重复。

\paragraph{实验协议。} 我们把 188 个样例分成五折，并保证同一条轨迹的所有样例落在同一折内，从而没有任何轨迹同时出现在训练与测试中。我们在四折上训练、给第五折打分，再把五折的打分结果汇总；下文每一个交叉验证数字都来自这个汇总。对外部集，我们在全部 188 个样例上训练一次、评测一次。\S\ref{sec:index} 说明哪些 writer 的数字不是交叉验证得到的。
